\documentclass[aspectratio=169,9pt]{beamer}

\usetheme{default}
\usefonttheme{professionalfonts}
\setbeamertemplate{navigation symbols}{}
\setbeamertemplate{footline}{\hfill\scriptsize\insertframenumber/\inserttotalframenumber\hspace{2mm}\vspace{1mm}}

\usepackage{amsmath,amssymb,bm,mathtools,array,booktabs}
\usepackage{xcolor}
\usepackage[most]{tcolorbox}

\definecolor{deepblue}{RGB}{0,70,150}
\definecolor{softblue}{RGB}{232,242,255}
\definecolor{softgreen}{RGB}{235,250,238}
\definecolor{softorange}{RGB}{255,245,230}
\definecolor{softpurple}{RGB}{245,238,255}
\definecolor{softgray}{RGB}{246,247,249}
\definecolor{darkred}{RGB}{160,35,25}
\definecolor{darkgreen}{RGB}{20,120,60}
\definecolor{darkpurple}{RGB}{95,45,150}
\definecolor{darkgray}{RGB}{80,80,80}

\setbeamerfont{frametitle}{series=\bfseries,size=\Large}
\setbeamercolor{frametitle}{fg=black}
\setbeamertemplate{itemize item}{\color{deepblue}$\blacktriangleright$}
\setbeamertemplate{itemize subitem}{\color{darkgreen}$\bullet$}
\setlength{\abovedisplayskip}{2pt}
\setlength{\belowdisplayskip}{2pt}
\setlength{\jot}{1pt}
\setlength{\tabcolsep}{3pt}

\newtcolorbox{bluebox}[1][]{colback=softblue,colframe=deepblue,boxrule=0.7pt,arc=1.3mm,left=1.4mm,right=1.4mm,top=0.8mm,bottom=0.8mm,#1}
\newtcolorbox{greenbox}[1][]{colback=softgreen,colframe=darkgreen,boxrule=0.7pt,arc=1.3mm,left=1.4mm,right=1.4mm,top=0.8mm,bottom=0.8mm,#1}
\newtcolorbox{orangebox}[1][]{colback=softorange,colframe=darkred,boxrule=0.7pt,arc=1.3mm,left=1.4mm,right=1.4mm,top=0.8mm,bottom=0.8mm,#1}
\newtcolorbox{purplebox}[1][]{colback=softpurple,colframe=darkpurple,boxrule=0.7pt,arc=1.3mm,left=1.4mm,right=1.4mm,top=0.8mm,bottom=0.8mm,#1}
\newtcolorbox{graybox}[1][]{colback=softgray,colframe=darkgray,boxrule=0.6pt,arc=1.3mm,left=1.4mm,right=1.4mm,top=0.8mm,bottom=0.8mm,#1}

\newcommand{\dd}{\,\mathrm{d}}
\newcommand{\grad}{\nabla}
\newcommand{\divg}{\nabla\!\cdot}
\newcommand{\E}{\bm{E}}
\newcommand{\J}{\bm{J}}
\newcommand{\q}{q}
\newcommand{\epsSi}{\epsilon_{\mathrm{Si}}}
\newcommand{\rhoDef}{\rho_{\mathrm{def}}}
\newcommand{\Omegae}{\Omega_e}
\newcommand{\vphi}{\phi}
\newcommand{\vect}[1]{\bm{#1}}
\newcommand{\mat}[1]{\mathbf{#1}}
\newcommand{\N}{\mat{N}}
\newcommand{\Be}{\mat{B}_e}
\newcommand{\Ke}{\mat{K}_e}
\newcommand{\Me}{\mat{M}_e}
\newcommand{\Ree}{\mat{R}_e}
\newcommand{\Ce}{\mat{C}_e}

\begin{document}

% ---------------- Module 2 ----------------
\begin{frame}{Module 2: Electrostatic bridge from radiation defects to field}
\footnotesize
\vspace{-3mm}
\begin{bluebox}
\centering
\textbf{Purpose:} convert mobile carriers, ionized dopants, and charged radiation-induced defects into a self-consistent electrostatic potential and electric field in silicon.
\end{bluebox}

\vspace{1mm}
\begin{columns}[T,totalwidth=\textwidth]
\begin{column}{0.50\textwidth}
\textbf{Space-charge density}
\begin{equation*}
\boxed{\rho = q(p-n+N_D^+-N_A^-) + \rho_{\mathrm{def}}}
\end{equation*}
\begin{itemize}\setlength\itemsep{0.4mm}
\item \(p\): mobile hole concentration.
\item \(n\): mobile electron concentration.
\item \(N_D^+\): ionized donor concentration.
\item \(N_A^-\): ionized acceptor concentration.
\item \(\rho_{\mathrm{def}}\): charge from radiation-induced defect populations.
\end{itemize}
\end{column}
\begin{column}{0.48\textwidth}
\textbf{Electrostatic output}
\begin{equation*}
\boxed{\rho(x,y,t)\longrightarrow \phi(x,y,t)\longrightarrow \E(x,y,t)}
\end{equation*}
\begin{equation*}
\boxed{\E=-\grad \phi}
\end{equation*}
\small
The module does not evolve defects by itself.  It receives defect populations from Module 3, converts them into charge density, then solves the field problem.
\end{column}
\end{columns}

\vspace{1mm}
\begin{orangebox}
\centering\scriptsize
\textbf{Core message:} radiation damage changes material charge; charge bends potential; potential gradients create the electric field that controls drift in Module 5.
\end{orangebox}
\end{frame}

\begin{frame}{Module 2: Electrostatic limit and Poisson equation}
\footnotesize
\vspace{-3mm}
\begin{bluebox}
\centering
Module 2 uses electroquasistatics: the field is charge-dominated, not induction-dominated.
\end{bluebox}

\vspace{1mm}
\begin{columns}[T,totalwidth=\textwidth]
\begin{column}{0.50\textwidth}
\textbf{Maxwell starting point}
\begin{align*}
\divg \vect{D} &= \rho, \\
\vect{D} &= \epsilon \E, \\
\E &= -\grad \phi.
\end{align*}
Substitution gives the electrostatic field equation:
\begin{equation*}
\boxed{\divg(\epsilon \grad \phi) = -\rho}
\end{equation*}
For uniform silicon permittivity,
\begin{equation*}
\boxed{\epsSi \grad^2\phi=-\rho.}
\end{equation*}
\end{column}
\begin{column}{0.48\textwidth}
\textbf{Why induction is neglected}
\begin{equation*}
\boxed{\frac{E_{\mathrm{ind}}}{E_{\mathrm{es}}}\sim
\left(\frac{L}{c_{\mathrm{Si}}\tau}\right)^2}
\end{equation*}
\begin{itemize}\setlength\itemsep{0.4mm}
\item \(L\): device length scale.
\item \(\tau\): time scale for charge redistribution.
\item \(c_{\mathrm{Si}}=1/\sqrt{\mu\epsilon}\): EM wave speed in silicon.
\end{itemize}
\small
If \(\tau \gg L/c_{\mathrm{Si}}\), the field adjusts much faster than the material charge evolves; Poisson's equation is the correct reduced model.
\end{column}
\end{columns}

\vspace{1mm}
\begin{greenbox}
\centering\scriptsize
\textbf{Field solver equation:}\quad
\(\displaystyle \frac{\partial^2\phi}{\partial x^2}+\frac{\partial^2\phi}{\partial y^2}= -\rho/\epsSi\).
\end{greenbox}
\end{frame}

\begin{frame}{Module 2: Defect species and effective charge state}
\footnotesize
\vspace{-3mm}
\begin{bluebox}
\centering
A defect species is a radiation-induced material-defect population, not a mobile carrier or dopant.
\end{bluebox}

\vspace{1mm}
\begin{columns}[T,totalwidth=\textwidth]
\begin{column}{0.52\textwidth}
\textbf{Defect-charge map}
\begin{equation*}
\boxed{\rho_{\mathrm{def}} = q\sum_{i=1}^{M} z_i C_i}
\end{equation*}
\small
Examples of \(C_i\):
\begin{center}
\scriptsize
\begin{tabular}{lll}
\toprule
Field & Meaning & Charge factor \\
\midrule
\(C_V\) & vacancy concentration & \(z_V\) \\
\(C_I\) & self-interstitial concentration & \(z_I\) \\
\(C_{V_2}\) & divacancy concentration & \(z_{V_2}\) \\
\(C_{VO}\) & vacancy-oxygen complex & \(z_{VO}\) \\
\bottomrule
\end{tabular}
\end{center}
\small
These are distinct from \(p,n,N_D^+,N_A^-\), which already appear separately in \(\rho\).
\end{column}
\begin{column}{0.46\textwidth}
\textbf{Why a vacancy is not automatically \(+q\)}
\small
A vacancy is a missing lattice atom, not a missing valence-band electron.  Removing a silicon atom changes a covalent bonding network: dangling bonds, lattice relaxation, and localized defect states may appear.
\begin{equation*}
\boxed{V^+,\quad V^0,\quad V^-,\quad V^{2-},\ldots}
\end{equation*}
\small
The dimensionless factor \(z_i\) represents the net electronic charge state after defect-state occupancy, not simply the fact that an atom is missing.
\end{column}
\end{columns}

\vspace{1mm}
\begin{purplebox}
\centering\scriptsize
\(z_i>0\): donor-like contribution; \(z_i<0\): acceptor-like contribution; \(z_i=0\): neutral in the reduced model.
\end{purplebox}
\end{frame}

\begin{frame}{Module 2: FEM weak form and matrix equation}
\footnotesize
\vspace{-3mm}
\begin{bluebox}
\centering
The continuous Poisson equation becomes a finite-dimensional system for nodal potentials.
\end{bluebox}

\vspace{1mm}
\begin{columns}[T,totalwidth=\textwidth]
\begin{column}{0.50\textwidth}
\textbf{Weak form}
\begin{equation*}
\divg(\epsSi \grad\phi)=-\rho
\end{equation*}
Multiply by test function \(w\), integrate by parts:
\begin{equation*}
\boxed{\int_{\Omega}\epsSi\grad w\cdot\grad\phi\,\dd\Omega
=\int_{\Omega}w\rho\,\dd\Omega + \int_{\Gamma_N}w\bar{d}\,\dd\Gamma}
\end{equation*}
where \(\bar{d}=\vect{n}\cdot\epsSi\grad\phi\) is prescribed normal displacement-flux data.
\end{column}
\begin{column}{0.48\textwidth}
\textbf{Linear triangular element}
\begin{equation*}
\phi_h(x,y)=\sum_{a=1}^{3}N_a(x,y)\Phi_a,
\qquad
w_h=N_b.
\end{equation*}
\begin{equation*}
\boxed{K^e_{ab}=\int_{\Omega_e}\epsSi\grad N_a\cdot\grad N_b\,\dd\Omega}
\end{equation*}
\begin{equation*}
\boxed{b^e_a=\int_{\Omega_e}N_a\rho\,\dd\Omega+\int_{\Gamma^e_N}N_a\bar{d}\,\dd\Gamma}
\end{equation*}
\begin{equation*}
\boxed{\mat{K}\vect{\Phi}=\vect{b}}
\end{equation*}
\end{column}
\end{columns}

\vspace{1mm}
\begin{orangebox}
\centering\scriptsize
Dirichlet contacts prescribe \(\phi\); Neumann/symmetry boundaries prescribe normal field.  Postprocess \(\E=-\grad\phi\) elementwise or at nodes.
\end{orangebox}
\end{frame}

% ---------------- Module 3 ----------------
\begin{frame}{Module 3: Defect evolution equation and source terms}
\footnotesize
\vspace{-3mm}
\begin{bluebox}
\centering
\textbf{Purpose:} evolve radiation-induced defect populations from local generation, diffusion, reaction, clustering, and annealing.
\end{bluebox}

\vspace{1mm}
\begin{columns}[T,totalwidth=\textwidth]
\begin{column}{0.52\textwidth}
\textbf{Single-species reduced model}
\begin{equation*}
\boxed{\frac{\partial C}{\partial t}
=\divg(D\grad C)+G-k_a C-R(C)}
\end{equation*}
\begin{itemize}\setlength\itemsep{0.35mm}
\item \(C(x,y,t)\): defect concentration.
\item \(D\): effective diffusivity.
\item \(G\): radiation-induced generation source.
\item \(k_a C\): first-order annealing/removal.
\item \(R(C)\): nonlinear reaction or clustering term.
\end{itemize}
\end{column}
\begin{column}{0.46\textwidth}
\textbf{Multi-species form}
\begin{equation*}
\boxed{\frac{\partial C_i}{\partial t}
=\divg(D_i\grad C_i)+G_i+\sum_r \nu_{ir}\mathcal{R}_r(\{C_j\})-k_i C_i}
\end{equation*}
\small
A reaction \(r\) may represent recombination, complex formation, dissociation, or trap creation.
\begin{equation*}
V+I\rightarrow \emptyset,
\qquad
V+O\rightarrow VO,
\qquad
V+V\rightarrow V_2.
\end{equation*}
\end{column}
\end{columns}

\vspace{1mm}
\begin{greenbox}
\centering\scriptsize
\textbf{Module handoff:} Module 3 outputs \(C_i(x,y,t)\); Module 2 maps them to \(\rho_{\mathrm{def}}=q\sum_i z_iC_i\).
\end{greenbox}
\end{frame}

\begin{frame}{Module 3: Effective transport and kinetic coefficients}
\footnotesize
\vspace{-3mm}
\begin{bluebox}
\centering
The defect PDE is simple in form, but its coefficients carry the material physics.
\end{bluebox}

\vspace{1mm}
\begin{columns}[T,totalwidth=\textwidth]
\begin{column}{0.50\textwidth}
\textbf{Thermally activated motion}
\begin{equation*}
\boxed{D_i(T)=D_{i0}\exp\!\left(-\frac{E_{m,i}}{k_B T}\right)}
\end{equation*}
\begin{equation*}
\boxed{k_i(T)=\nu_i\exp\!\left(-\frac{E_{a,i}}{k_B T}\right)}
\end{equation*}
\small
\(E_{m,i}\) is a migration barrier; \(E_{a,i}\) is an annealing or reaction activation barrier.
\end{column}
\begin{column}{0.48\textwidth}
\textbf{Possible field dependence}
\small
If charged-defect drift is included,
\begin{equation*}
\boxed{\J_i=-D_i\grad C_i+\mu_i z_i q C_i\E}
\end{equation*}
\begin{equation*}
\boxed{\frac{\partial C_i}{\partial t}=-\divg\J_i+G_i+\text{reactions}}
\end{equation*}
Einstein relation, when applicable:
\begin{equation*}
\boxed{\mu_i=\frac{D_i}{k_B T}}
\end{equation*}
\small
The reduced baseline may omit this drift term; the coupled model can feed \(\E\) from Module 2 back into Module 3.
\end{column}
\end{columns}

\vspace{1mm}
\begin{purplebox}
\centering\scriptsize
\textbf{Interpretation:} Module 3 is where radiation damage becomes an evolving material state, not merely a static fluence-to-defect conversion.
\end{purplebox}
\end{frame}

\begin{frame}{Module 3: FEM weak form for diffusion--reaction defects}
\footnotesize
\vspace{-3mm}
\begin{bluebox}
\centering
The FEM form converts a defect transport PDE into mass, diffusion, reaction, and source matrices.
\end{bluebox}

\vspace{1mm}
\begin{columns}[T,totalwidth=\textwidth]
\begin{column}{0.50\textwidth}
\textbf{Weak form}
\begin{equation*}
\frac{\partial C}{\partial t}=\divg(D\grad C)+G-k_aC
\end{equation*}
Multiply by \(w\), integrate by parts:
\begin{equation*}
\boxed{\int_\Omega w\frac{\partial C}{\partial t}\dd\Omega
+\int_\Omega D\grad w\cdot\grad C\dd\Omega
+\int_\Omega wk_aC\dd\Omega
=\int_\Omega wG\dd\Omega}
\end{equation*}
Natural zero-flux boundary:
\begin{equation*}
\boxed{\vect{n}\cdot D\grad C=0}
\end{equation*}
\end{column}
\begin{column}{0.48\textwidth}
\textbf{Element matrices}
\begin{equation*}
C_h=\sum_a N_a C_a
\end{equation*}
{\scriptsize
\begin{align*}
M^e_{ab}&=\int_{\Omega_e}N_aN_b\dd\Omega,\\
K^e_{ab}&=\int_{\Omega_e}D\grad N_a\cdot\grad N_b\dd\Omega,\\
R^e_{ab}&=\int_{\Omega_e}k_a N_aN_b\dd\Omega,\\
g^e_a&=\int_{\Omega_e}N_aG\dd\Omega.
\end{align*}}
\begin{equation*}
\boxed{\mat{M}\dot{\vect{C}}+(\mat{K}+\mat{R})\vect{C}=\vect{g}}
\end{equation*}
\end{column}
\end{columns}
\end{frame}

\begin{frame}{Module 3: Time stepping, verification, and coupling}
\footnotesize
\vspace{-3mm}
\begin{bluebox}
\centering
Backward Euler gives a stable update for diffusion and annealing over stiff time scales.
\end{bluebox}

\vspace{1mm}
\begin{columns}[T,totalwidth=\textwidth]
\begin{column}{0.52\textwidth}
\textbf{Implicit defect update}
\begin{equation*}
\boxed{\left[\mat{M}+\Delta t(\mat{K}+\mat{R})\right]\vect{C}^{n+1}
=\mat{M}\vect{C}^{n}+\Delta t\,\vect{g}^{n+1}}
\end{equation*}
\small
For nonlinear reactions,
\begin{equation*}
\boxed{\mat{F}(\vect{C}^{n+1})=\vect{0}}
\end{equation*}
may be solved by fixed-point iteration, Newton iteration, or operator splitting.
\end{column}
\begin{column}{0.46\textwidth}
\textbf{Verification logic}
\begin{itemize}\setlength\itemsep{0.45mm}
\item \textbf{Pure annealing:} with \(D=0, G=0\), recover \(C(t)=C_0e^{-k_at}\).
\item \textbf{Uniform preservation:} with zero flux and no source, uniform \(C\) remains uniform.
\item \textbf{Diffusion inventory:} with zero flux and no reactions, \(\int_\Omega C\dd\Omega\) is conserved.
\item \textbf{Coupling sanity:} increasing \(C_i\) should increase \(\rho_{\mathrm{def}}\) in Module 2 according to \(qz_iC_i\).
\end{itemize}
\end{column}
\end{columns}

\vspace{1mm}
\begin{orangebox}
\centering\scriptsize
Module 3 supplies the evolving damage fields that drive electrostatics, thermal coefficients, and recombination/trapping in later modules.
\end{orangebox}
\end{frame}

% ---------------- Module 4 ----------------
\begin{frame}{Module 4: Thermal response and ballistic--diffusive motivation}
\footnotesize
\vspace{-3mm}
\begin{bluebox}
\centering
\textbf{Purpose:} convert deposited and thermalized radiation energy into an evolving temperature field, while retaining quasi-ballistic corrections when local Fourier transport is insufficient.
\end{bluebox}

\vspace{1mm}
\begin{columns}[T,totalwidth=\textwidth]
\begin{column}{0.50\textwidth}
\textbf{Continuum energy balance}
\begin{equation*}
\boxed{C_T\frac{\partial T}{\partial t}+\divg\vect{q}=Q_T}
\end{equation*}
Fourier closure:
\begin{equation*}
\boxed{\vect{q}=-k\grad T}
\end{equation*}
so
\begin{equation*}
\boxed{C_T\frac{\partial T}{\partial t}=\divg(k\grad T)+Q_T.}
\end{equation*}
\end{column}
\begin{column}{0.48\textwidth}
\textbf{When Fourier is not enough}
\begin{equation*}
\boxed{\mathrm{Kn}=\frac{\lambda}{L}}
\end{equation*}
\begin{itemize}\setlength\itemsep{0.35mm}
\item \(\lambda\): heat-carrier mean free path.
\item \(L\): thermal length scale or device dimension.
\item \(\mathrm{Kn}\ll1\): diffusive limit.
\item \(\mathrm{Kn}\sim0.1-1\): ballistic memory matters.
\end{itemize}
\small
Radiation energy can initially reside in nonlocal carriers before relaxing into the thermal bath.
\end{column}
\end{columns}

\vspace{1mm}
\begin{greenbox}
\centering\scriptsize
Module 4 outputs \(T(x,y,t)\), which controls defect kinetics, mobilities, recombination rates, and coupled thermal feedback.
\end{greenbox}
\end{frame}

\begin{frame}{Module 4: Reduced ballistic--diffusive source structure}
\footnotesize
\vspace{-3mm}
\begin{bluebox}
\centering
The reduced model keeps a heat-equation solve but modifies the source with ballistic energy exchange.
\end{bluebox}

\vspace{1mm}
\begin{columns}[T,totalwidth=\textwidth]
\begin{column}{0.50\textwidth}
\textbf{Energy and flux split}
\begin{equation*}
\boxed{u=u_b+u_m},\qquad
\boxed{\vect{q}=\vect{q}_b+\vect{q}_m}
\end{equation*}
\begin{equation*}
\boxed{u_m=C_T T_m},\qquad
\boxed{\vect{q}_m=-k\grad T_m}
\end{equation*}
\small
The subscript \(b\) denotes ballistic/nonlocal carrier energy; \(m\) denotes medium thermal energy.
\end{column}
\begin{column}{0.48\textwidth}
\textbf{Effective medium-temperature PDE}
\begin{equation*}
\boxed{C_T\frac{\partial T_m}{\partial t}
=\divg(k\grad T_m)
+Q_{\mathrm{th}}
-\divg\vect{q}_b
-\frac{\partial u_b}{\partial t}}
\end{equation*}
\small
The ballistic correction can be treated as an effective source:
\begin{equation*}
\boxed{Q_{\mathrm{eff}}=Q_{\mathrm{th}}-\divg\vect{q}_b-\partial_t u_b}
\end{equation*}
\begin{equation*}
\boxed{C_T\partial_t T_m=\divg(k\grad T_m)+Q_{\mathrm{eff}}}
\end{equation*}
\end{column}
\end{columns}

\vspace{1mm}
\begin{purplebox}
\centering\scriptsize
The FEM implementation can solve the same parabolic structure as Fourier heat conduction, while the source term stores the ballistic physics.
\end{purplebox}
\end{frame}

\begin{frame}{Module 4: FEM weak form and thermal matrices}
\footnotesize
\vspace{-3mm}
\begin{bluebox}
\centering
Temperature is approximated on the shared 2D mesh using the same finite-element machinery as the other field modules.
\end{bluebox}

\vspace{1mm}
\begin{columns}[T,totalwidth=\textwidth]
\begin{column}{0.52\textwidth}
\textbf{Weak form}
\begin{equation*}
C_T\frac{\partial T}{\partial t}=\divg(k\grad T)+Q_{\mathrm{eff}}
\end{equation*}
\begin{equation*}
\boxed{\int_\Omega w C_T\frac{\partial T}{\partial t}\dd\Omega
+\int_\Omega k\grad w\cdot\grad T\dd\Omega
=\int_\Omega wQ_{\mathrm{eff}}\dd\Omega+\int_{\Gamma_N}w\bar{q}\dd\Gamma}
\end{equation*}
\small
Natural insulated boundary:
\begin{equation*}
\boxed{\vect{n}\cdot k\grad T=0.}
\end{equation*}
\end{column}
\begin{column}{0.46\textwidth}
\textbf{Element arrays}
\begin{equation*}
T_h=\sum_aN_aT_a
\end{equation*}
{\scriptsize
\begin{align*}
M^e_{ab}&=\int_{\Omega_e}C_TN_aN_b\dd\Omega,\\
K^e_{ab}&=\int_{\Omega_e}k\grad N_a\cdot\grad N_b\dd\Omega,\\
f^e_a&=\int_{\Omega_e}N_aQ_{\mathrm{eff}}\dd\Omega.
\end{align*}}
\begin{equation*}
\boxed{\mat{M}_T\dot{\vect{T}}+\mat{K}_T\vect{T}=\vect{f}_T}
\end{equation*}
\end{column}
\end{columns}
\end{frame}

\begin{frame}{Module 4: Thermal time stepping and checks}
\footnotesize
\vspace{-3mm}
\begin{bluebox}
\centering
The thermal solve advances temperature with sources from radiation thermalization and ballistic energy exchange.
\end{bluebox}

\vspace{1mm}
\begin{columns}[T,totalwidth=\textwidth]
\begin{column}{0.52\textwidth}
\textbf{Backward-Euler update}
\begin{equation*}
\boxed{\left(\mat{M}_T+\Delta t\mat{K}_T\right)\vect{T}^{n+1}
=\mat{M}_T\vect{T}^{n}+\Delta t\,\vect{f}_T^{n+1}}
\end{equation*}
\textbf{Source construction}
\begin{equation*}
\boxed{Q_{\mathrm{th}}=\eta_{\mathrm{th}}\dot{q}_{\mathrm{rad}}^{2D}+Q_{\mathrm{device}}}
\end{equation*}
\begin{equation*}
\boxed{Q_{\mathrm{eff}}=Q_{\mathrm{th}}+Q_{\mathrm{ballistic}}}
\end{equation*}
\end{column}
\begin{column}{0.46\textwidth}
\textbf{Verification logic}
\begin{itemize}\setlength\itemsep{0.45mm}
\item \textbf{Uniform equilibrium:} no source and insulated boundaries preserve uniform \(T\).
\item \textbf{Fourier Gaussian diffusion:} in the ballistic-free limit, spreading follows diffusion scaling.
\item \textbf{Uniform source:} spatially uniform \(Q\) causes uniform temperature rise.
\item \textbf{Energy check:} integrated temperature rise matches input energy when heat loss is absent.
\end{itemize}
\end{column}
\end{columns}

\vspace{1mm}
\begin{orangebox}
\centering\scriptsize
\textbf{Coupling:} \(T\) modifies defect diffusivity/annealing in Module 3 and carrier mobility/recombination in Module 5.
\end{orangebox}
\end{frame}

% ---------------- Module 5 ----------------
\begin{frame}{Module 5: Carrier drift--diffusion with radiation defects}
\footnotesize
\vspace{-3mm}
\begin{bluebox}
\centering
\textbf{Purpose:} evolve mobile electrons and holes under diffusion, electric-field drift, generation, and defect-assisted recombination/trapping.
\end{bluebox}

\vspace{1mm}
\begin{columns}[T,totalwidth=\textwidth]
\begin{column}{0.50\textwidth}
\textbf{Continuity equations}
\begin{align*}
\boxed{\frac{\partial n}{\partial t}}&\boxed{=\frac{1}{q}\divg\J_n+G_n-R_n},\\
\boxed{\frac{\partial p}{\partial t}}&\boxed{=-\frac{1}{q}\divg\J_p+G_p-R_p.}
\end{align*}
\small
The sign difference follows from electron negative charge and hole positive charge.
\end{column}
\begin{column}{0.48\textwidth}
\textbf{Current densities}
\begin{align*}
\boxed{\J_n}&\boxed{=q\mu_n n\E+qD_n\grad n},\\
\boxed{\J_p}&\boxed{=q\mu_p p\E-qD_p\grad p.}
\end{align*}
\begin{equation*}
\boxed{\E=-\grad\phi}
\end{equation*}
\small
Module 2 supplies \(\phi\) and \(\E\); Module 3 supplies defect densities that influence recombination/trapping.
\end{column}
\end{columns}

\vspace{1mm}
\begin{greenbox}
\centering\scriptsize
Carrier transport is where radiation-induced field distortion and defect-assisted recombination become electrical degradation observables.
\end{greenbox}
\end{frame}

\begin{frame}{Module 5: Recombination, generation, and defect coupling}
\footnotesize
\vspace{-3mm}
\begin{bluebox}
\centering
Defects affect mobile carriers both electrostatically and kinetically.
\end{bluebox}

\vspace{1mm}
\begin{columns}[T,totalwidth=\textwidth]
\begin{column}{0.50\textwidth}
\textbf{Reduced recombination/trapping model}
\begin{equation*}
\boxed{R_n\approx \frac{n-n_{\mathrm{eq}}}{\tau_n(C_i,T)}},\qquad
\boxed{R_p\approx \frac{p-p_{\mathrm{eq}}}{\tau_p(C_i,T)}}
\end{equation*}
\small
Defect concentrations can shorten lifetimes:
\begin{equation*}
\boxed{\frac{1}{\tau_n}=\frac{1}{\tau_{n0}}+\sum_i \alpha_{n,i}C_i}
\end{equation*}
\begin{equation*}
\boxed{\frac{1}{\tau_p}=\frac{1}{\tau_{p0}}+\sum_i \alpha_{p,i}C_i.}
\end{equation*}
\end{column}
\begin{column}{0.48\textwidth}
\textbf{Two damage pathways}
\begin{itemize}\setlength\itemsep{0.55mm}
\item \textbf{Field pathway:} \(C_i\rightarrow\rho_{\mathrm{def}}\rightarrow \phi,\E\rightarrow\) drift.
\item \textbf{Lifetime pathway:} \(C_i\rightarrow\tau_n,\tau_p\rightarrow R_n,R_p\rightarrow\) carrier loss.
\item \textbf{Thermal pathway:} \(T\rightarrow\mu_n,\mu_p,D_n,D_p,\tau_n,\tau_p\).
\end{itemize}
\begin{equation*}
\boxed{G_{n,p}=G_{\mathrm{rad}}+G_{\mathrm{device}}}
\end{equation*}
\small
Generation may include radiation-induced carriers, optical excitation, or device injection depending on the case.
\end{column}
\end{columns}
\end{frame}

\begin{frame}{Module 5: FEM weak form and carrier matrices}
\footnotesize
\vspace{-3mm}
\begin{bluebox}
\centering
The drift--diffusion PDE gives a mass matrix, diffusion matrix, drift/advection matrix, recombination matrix, and source vector.
\end{bluebox}

\vspace{1mm}
\begin{columns}[T,totalwidth=\textwidth]
\begin{column}{0.53\textwidth}
\textbf{Electron weak form, schematic}
\begin{equation*}
\frac{\partial n}{\partial t}=\divg(D_n\grad n+\mu_nn\E)+G_n-R_n
\end{equation*}
{\scriptsize
\begin{align*}
\int_\Omega w\partial_t n\dd\Omega
+\int_\Omega D_n\grad w\cdot\grad n\dd\Omega&\nonumber\\
-\int_\Omega \mu_n n\E\cdot\grad w\dd\Omega
&=\int_\Omega w(G_n-R_n)\dd\Omega .
\end{align*}}
\small
Boundary terms encode contact injection, recombination surfaces, or insulating zero-flux edges.
\end{column}
\begin{column}{0.45\textwidth}
\textbf{Discrete form}
\begin{equation*}
\boxed{\mat{M}_n\dot{\vect{n}}+(\mat{K}_n+\mat{A}_n+\mat{R}_n)\vect{n}=\vect{g}_n}
\end{equation*}
{\scriptsize
\begin{align*}
M^e_{ab}&=\int_{\Omega_e}N_aN_b\dd\Omega,\\
K^e_{ab}&=\int_{\Omega_e}D_n\grad N_a\cdot\grad N_b\dd\Omega,\\
A^e_{ab}&=-\int_{\Omega_e}\mu_n N_b\E\cdot\grad N_a\dd\Omega.
\end{align*}}
\small
A parallel system is assembled for holes with the proper drift sign.
\end{column}
\end{columns}
\end{frame}

\begin{frame}{Module 5: Time stepping, currents, and verification}
\footnotesize
\vspace{-3mm}
\begin{bluebox}
\centering
The transport solve advances carrier density and then postprocesses current density and degradation metrics.
\end{bluebox}

\vspace{1mm}
\begin{columns}[T,totalwidth=\textwidth]
\begin{column}{0.53\textwidth}
\textbf{Backward-Euler electron update}
\begin{equation*}
\boxed{\left[\mat{M}_n+\Delta t(\mat{K}_n+\mat{A}_n+\mat{R}_n)\right]\vect{n}^{n+1}
=\mat{M}_n\vect{n}^{n}+\Delta t\vect{g}_n^{n+1}}
\end{equation*}
\small
Analogous update for \(\vect{p}^{n+1}\).

\vspace{1mm}
\textbf{Current postprocessing}
\begin{align*}
\boxed{\J_n}&\boxed{=q\mu_n n\E+qD_n\grad n},\\
\boxed{\J_p}&\boxed{=q\mu_p p\E-qD_p\grad p.}
\end{align*}
\end{column}
\begin{column}{0.45\textwidth}
\textbf{Verification logic}
\begin{itemize}\setlength\itemsep{0.4mm}
\item No field, no source, no recombination: uniform \(n,p\) remain constant.
\item Pure recombination: recover exponential lifetime decay.
\item Pure diffusion: carrier inventory conserved under zero-flux boundaries.
\item Uniform field: drift/current sign matches electron and hole charge conventions.
\end{itemize}
\end{column}
\end{columns}

\vspace{1mm}
\begin{orangebox}
\centering\scriptsize
Module 5 turns fields and material damage into measurable electrical outputs: current, carrier loss, leakage trends, and charge-collection degradation.
\end{orangebox}
\end{frame}

% ---------------- Module 6 ----------------
\begin{frame}{Module 6: Coupled multiphysics state vector}
\footnotesize
\vspace{-3mm}
\begin{bluebox}
\centering
\textbf{Purpose:} integrate electrostatics, defect evolution, thermal transport, and carrier drift--diffusion on a shared 2D FEM mesh.
\end{bluebox}

\vspace{1mm}
\begin{columns}[T,totalwidth=\textwidth]
\begin{column}{0.50\textwidth}
\textbf{Coupled unknown fields}
\begin{equation*}
\boxed{\vect{u}=\left[\vect{C}_1,\ldots,\vect{C}_M,\vect{\Phi},\vect{T},\vect{n},\vect{p}\right]^T}
\end{equation*}
\begin{itemize}\setlength\itemsep{0.4mm}
\item \(\vect{C}_i\): defect concentrations.
\item \(\vect{\Phi}\): electrostatic potential.
\item \(\vect{T}\): temperature.
\item \(\vect{n},\vect{p}\): carrier densities.
\end{itemize}
\end{column}
\begin{column}{0.48\textwidth}
\textbf{Shared-mesh philosophy}
\begin{equation*}
\boxed{\text{same nodes/elements} \Rightarrow \text{direct field transfer}}
\end{equation*}
\small
When all modules use the same triangular mesh, couplings such as
\begin{equation*}
C_i\rightarrow \rho_{\mathrm{def}},\quad
\phi\rightarrow \E,\quad
T\rightarrow D_i,k_i,\mu_n,\mu_p
\end{equation*}
are evaluated without interpolation errors between incompatible grids.
\end{column}
\end{columns}

\vspace{1mm}
\begin{greenbox}
\centering\scriptsize
Module 6 is not a new physics law; it is the numerical architecture that makes Modules 2--5 talk to each other consistently.
\end{greenbox}
\end{frame}

\begin{frame}{Module 6: Coupled block residual structure}
\footnotesize
\vspace{-3mm}
\begin{bluebox}
\centering
The fully coupled system can be written as interacting residual blocks.
\end{bluebox}

\vspace{1mm}
\begin{equation*}
\boxed{
\vect{R}(\vect{u})=
\begin{bmatrix}
\vect{R}_{C_1}\\[-1mm]
\vdots\\[-1mm]
\vect{R}_{C_M}\\[-1mm]
\vect{R}_{\phi}\\[-1mm]
\vect{R}_{T}\\[-1mm]
\vect{R}_{n}\\[-1mm]
\vect{R}_{p}
\end{bmatrix}
=\vect{0}}
\end{equation*}

\vspace{1mm}
\begin{columns}[T,totalwidth=\textwidth]
\begin{column}{0.49\textwidth}
\textbf{Representative blocks}
\begin{align*}
\vect{R}_{\phi}&=\mat{K}_{\phi}\vect{\Phi}-\vect{b}_{\rho}(\vect{C}_i,\vect{n},\vect{p}),\\
\vect{R}_{C_i}&=\mat{M}_C\dot{\vect{C}}_i+\mat{K}_{C_i}(T,\E)\vect{C}_i-\vect{g}_{C_i},\\
\vect{R}_{T}&=\mat{M}_T\dot{\vect{T}}+\mat{K}_T\vect{T}-\vect{f}_T.
\end{align*}
\end{column}
\begin{column}{0.49\textwidth}
\textbf{Carrier blocks}
\begin{align*}
\vect{R}_{n}&=\mat{M}_n\dot{\vect{n}}+\left(\mat{K}_n+\mat{A}_n(\E)+\mat{R}_n(C_i,T)\right)\vect{n}-\vect{g}_n,\\
\vect{R}_{p}&=\mat{M}_p\dot{\vect{p}}+\left(\mat{K}_p+\mat{A}_p(\E)+\mat{R}_p(C_i,T)\right)\vect{p}-\vect{g}_p.
\end{align*}
\end{column}
\end{columns}

\vspace{1mm}
\begin{purplebox}
\centering\scriptsize
The residual view supports either a monolithic nonlinear solve or a staggered operator-splitting solve.
\end{purplebox}
\end{frame}

\begin{frame}{Module 6: Staggered execution order in MATLAB}
\footnotesize
\vspace{-3mm}
\begin{bluebox}
\centering
The computer must choose an order even though the physical system is coupled.
\end{bluebox}

\vspace{1mm}
\begin{columns}[T,totalwidth=\textwidth]
\begin{column}{0.52\textwidth}
\textbf{One time step, reduced staggered loop}
\begin{align*}
\text{1. }&\rho^n=q(p^n-n^n+N_D^+-N_A^-)+q\sum_i z_iC_i^n,\\
\text{2. }&\mat{K}_{\phi}\vect{\Phi}^n=\vect{b}_{\rho}^n,
\qquad \E^n=-\grad\phi^n,\\
\text{3. }&\vect{C}_i^{n}\rightarrow\vect{C}_i^{n+1}\quad\text{using }T^n,\E^n,\\
\text{4. }&\vect{T}^{n}\rightarrow\vect{T}^{n+1}\quad\text{using heat sources},\\
\text{5. }&\vect{n}^{n},\vect{p}^{n}\rightarrow\vect{n}^{n+1},\vect{p}^{n+1}\quad\text{using }\E^n,C_i^{n+1},T^{n+1},\\
\text{6. }&\text{recompute }\rho^{n+1},\phi^{n+1},\E^{n+1}.
\end{align*}
\end{column}
\begin{column}{0.46\textwidth}
\textbf{Optional inner coupling iteration}
\begin{equation*}
\boxed{\vect{u}^{n+1,k}\rightarrow\vect{u}^{n+1,k+1}}
\end{equation*}
\small
Repeat the module sequence until
\begin{equation*}
\boxed{\frac{\|\vect{u}^{k+1}-\vect{u}^{k}\|}{\|\vect{u}^{k+1}\|}<\epsilon_{\mathrm{couple}}}
\end{equation*}
\small
This improves consistency when field-defect-carrier feedback is strong.
\end{column}
\end{columns}

\vspace{1mm}
\begin{orangebox}
\centering\scriptsize
\textbf{Key point:} defects reshape the field, and the updated field may feed back into charged-defect drift or kinetic coefficients in the next substep or iteration.
\end{orangebox}
\end{frame}

\begin{frame}{Module 6: Transfer policy, diagnostics, and summary}
\footnotesize
\vspace{-3mm}
\begin{bluebox}
\centering
A coupled solver is only credible if field transfers and conservation diagnostics are explicit.
\end{bluebox}

\vspace{1mm}
\begin{columns}[T,totalwidth=\textwidth]
\begin{column}{0.51\textwidth}
\textbf{Transfer and projection rules}
\begin{itemize}\setlength\itemsep{0.45mm}
\item Same mesh: pass nodal fields directly.
\item Different mesh: project by interpolation or conservative \(L^2\) projection.
\item Element values: evaluate gradients and fluxes by element shape-function derivatives.
\item Boundary data: map contact potentials, heat fluxes, and carrier fluxes consistently.
\end{itemize}
\end{column}
\begin{column}{0.47\textwidth}
\textbf{Diagnostics}
\begin{align*}
Q_{\mathrm{tot}}&=\int_{\Omega}\rho\dd\Omega,\\
N_{C_i}&=\int_{\Omega}C_i\dd\Omega,\\
U_T&=\int_{\Omega}C_TT\dd\Omega.
\end{align*}
\small
Track charge consistency, defect inventory, thermal energy, and carrier balance after each coupled step.
\end{column}
\end{columns}

\vspace{1mm}
\begin{greenbox}
\small
\textbf{Final summary:} Module 2 solves fields from charge; Module 3 evolves defect populations; Module 4 evolves temperature and nonlocal thermalization; Module 5 evolves carriers and current; Module 6 couples them on a shared FEM infrastructure.
\end{greenbox}
\end{frame}

\end{document}
